% 2018 全国Ⅲ卷文科/理科第3题榫卯结构与俯视图选项
\begin{tikzpicture}[line cap=round, line join=round, line width=0.6pt, font=\small]

% 上部：木构件榫卯示意立体图
\begin{scope}[shift={(0.8, 0)}]
  % 主体长方体
  % 前表面 (0,0) 到 (3.0, 1.3)
  \draw (0,0) rectangle (3.0,1.3);
  % 顶面与侧面可见棱
  \draw (0,1.3) -- (1.2,2.0) -- (4.2,2.0) -- (3.0,1.3);
  \draw (3.0,0) -- (4.2,0.7) -- (4.2,2.0);
  % 内部不可见虚线棱
  \draw[dashed] (0,0) -- (1.2,0.7) -- (4.2,0.7);
  \draw[dashed] (1.2,0.7) -- (1.2,2.0);

  % 榫头（右侧凸出小长方体）
  % 榫头在右表面 (x从3.0到4.2, y从0.7到2.0斜面)，向右延伸 dx=1.2
  % 根部前后缩进，高度 0.3 到 1.0 (主体高度为 0 到 1.3)
  % 根部四个顶点：
  % 根前下: (3.3, 0.45), 根前上: (3.3, 1.15)
  % 根后下: (3.9, 0.8),  根后上: (3.9, 1.5)
  % 凸出端面（向右平移 1.2）:
  % 端前下: (4.5, 0.45), 端前上: (4.5, 1.15)
  % 端后下: (5.1, 0.8),  端后上: (5.1, 1.5)
  
  % 凸出端面
  \draw (4.5,0.45) -- (4.5,1.15) -- (5.1,1.5) -- (5.1,0.8) -- cycle;
  % 榫头可见连接棱
  \draw (3.3,1.15) -- (4.5,1.15);
  \draw (3.9,1.5) -- (5.1,1.5);
  \draw (3.3,0.45) -- (4.5,0.45);
  % 榫头不可见虚线棱
  \draw[dashed] (3.3,0.45) -- (3.9,0.8) -- (5.1,0.8);
  \draw[dashed] (3.3,0.45) -- (3.3,1.15);
  \draw[dashed] (3.9,0.8) -- (3.9,1.5);

  % 俯视方向箭头与文字
  \draw[->, >=stealth, line width=0.7pt] (4.4,2.7) -- (4.4,1.75);
  \node[anchor=west, align=center] at (4.45,2.25) {\footnotesize 俯\\[-2pt]\footnotesize 视\\[-2pt]\footnotesize 方\\[-2pt]\footnotesize 向};
\end{scope}

% 下部：四个选项俯视图
% (A)
\begin{scope}[shift={(0, -2.2)}]
  \draw (0,0) rectangle (2.8,1.3);
  \draw[dashed] (0,0.95) -- (1.2,0.95) -- (1.2,0.35) -- (0,0.35);
  \node[left] at (-0.1, 0.65) {(A)};
\end{scope}

% (B)
\begin{scope}[shift={(4.0, -2.2)}]
  \draw (0,0) rectangle (2.8,1.3);
  \draw (0,0.95) -- (1.2,0.95) -- (1.2,0.35) -- (0,0.35);
  \node[left] at (-0.1, 0.65) {(B)};
\end{scope}

% (C)
\begin{scope}[shift={(0, -4.0)}]
  \draw (0,0) rectangle (2.8,1.3);
  \draw[dashed] (0,0.8) -- (1.2,0.8) -- (1.2,0);
  \node[left] at (-0.1, 0.65) {(C)};
\end{scope}

% (D)
\begin{scope}[shift={(4.0, -4.0)}]
  \draw (0,0) -- (2.8,0) -- (2.8,1.3) -- (0,1.3) -- (0,0.95) -- (1.2,0.95) -- (1.2,0.35) -- (0,0.35) -- cycle;
  \node[left] at (-0.1, 0.65) {(D)};
\end{scope}

\end{tikzpicture}
